\scanpage{216}
\mspage{1 - 48 -}\hfill n\textsuperscript{o} 413

17. \uline{Noyaux de groupe}.

DÉFINITION 5.- \uline{On appelle noyau de classe} $\mathrm C^{\mathrm r}$ \uline{une variété} $\mathrm V$ \uline{de classe} $\mathrm C^{\mathrm r}$ \uline{munie d'un point marqué} $\mathrm e$. \uline{Un noyau de groupe de classe} $\mathrm C^{\mathrm r}$ \uline{est un noyau} $\mathrm G$ \uline{muni d'une application} $\unsure{\vartheta|\theta}_{\mathrm G}$ \uline{de classe} $\mathrm C^{\mathrm r}$ \uline{de} $\mathrm G$ \uline{dans} $\mathrm G$ \uline{et d'une application} $\mathrm m_{\mathrm G}$ \uline{de classe} $\mathrm C^{\mathrm r}$ \uline{d'un ouvert} $\Omega(\mathrm G)$ \uline{de} $\mathrm G\times\mathrm G$ \uline{dans} $\mathrm G$ \uline{satisfaisant aux conditions suivantes} :

\begingroup
\makeatletter\tagsleft@true\makeatother
(i) $\Omega$ \uline{contient les points} $(\mathrm x,\vartheta(\mathrm x))$ \uline{et} $(\vartheta(\mathrm x),\mathrm x)$ \uline{et on a}
\[
\mathrm m(\mathrm x,\vartheta(\mathrm x))=\mathrm m(\vartheta(\mathrm x),\mathrm x)=\mathrm e \tag*{(36)}
\]
\uline{pour tout} $\mathrm x\in\mathrm G$ ;

(ii) $\Omega$ \uline{contient les points} $(\mathrm e,\mathrm x)$ \uline{et} $(\mathrm x,\mathrm e)$ \uline{et on a}
\margin{\handwritten{\unsure{$\parallel$|T|$\pi$}}}
\[
\mathrm m(\mathrm e,\mathrm x)=\mathrm m(\mathrm x,\mathrm e)=\mathrm x \tag*{(37)}
\]
\uline{pour tout} $\mathrm x\in\mathrm G$ ;

(iii) \uline{si} $\mathrm x,\mathrm y,\mathrm z$ \uline{sont trois éléments de} $\mathrm G$ \uline{tels que} $(\mathrm x,\mathrm y)$, $(\mathrm y,\mathrm z)$, $(\mathrm m(\mathrm x,\mathrm y),\mathrm z)$ \uline{et} $\mathrm m(\mathrm x,\mathrm m(\mathrm y,\mathrm z))$ \uline{appartiennent à} $\Omega$\handwritten{$\unsure{\searrow|,}$} \uline{on a} :
\[
\mathrm m(\mathrm m(\mathrm x,\mathrm y),\mathrm z)=\mathrm m(\mathrm x,\mathrm m(\mathrm y,\mathrm z)). \tag*{(38)}
\]
\endgroup

\[
\begin{aligned}
\text{On a }\mathrm e=\mathrm m(\mathrm e,\vartheta(\mathrm e))=\vartheta(\mathrm e)
\quad\text{et }\mathrm x&=\mathrm m(\mathrm e,\mathrm x)=\mathrm m(\mathrm m(\vartheta(\vartheta(\mathrm x)),\vartheta(\mathrm x)),\mathrm x)=\\
&=\mathrm m(\vartheta(\vartheta(\mathrm x)),\mathrm m(\vartheta(\mathrm x),\mathrm x)=\mathrm m(\vartheta(\vartheta(\mathrm x)),\mathrm e)=\vartheta(\vartheta(\mathrm x)).
\end{aligned}
\]

\uline{Exemple} 1.- Une variété de groupe est un noyau de groupe.

\uline{Exemple} 2.- Soient $\mathrm X$ une variété, $\mathrm e\in\mathrm X$ et $\mathrm m$ une application d'un ouvert $\Omega$ de $\mathrm X\times\mathrm X$ dans $\mathrm X$ satisfaisant aux conditions (ii) et (iii) de la définition 5. Le lemme 1 du n\textsuperscript{o} 3 montre qu'il existe un voisinage ouvert $\mathrm G$ de $\mathrm e$ dans $\mathrm X$ et un automorphisme $\vartheta$ de $\mathrm G$ sur lui-même satisfaisant à (i). Il suffit alors de prendre pour $\Omega(\mathrm G)$ l'ensemble des $(\mathrm x,\mathrm y)\in(\mathrm G\times\mathrm G)\cap\Omega$ tels que $\mathrm m(\mathrm x,\mathrm y)\in\mathrm G$ et pour $\mathrm m_{\mathrm G}$ la restriction de $\mathrm m$ à $\Omega(\mathrm G)$ pour obtenir sur $\mathrm G$ une structure de noyau de groupe.
% archive sheet number: 215

